\documentclass[../main]{subfiles}
\begin{document}

\chapter{Leyes Electromagnéticas en Máquinas Eléctricas}

\section{Teoría Electromagnética Aplicada a Máquinas Eléctricas}

Las leyes electromagnéticas son fundamentales para entender el funcionamiento de las máquinas eléctricas. Estas leyes describen cómo se comportan los campos eléctricos y magnéticos, y cómo interactúan con los materiales conductores y magnéticos. En este capítulo, estudiaremos las principales leyes que rigen el comportamiento de las máquinas eléctricas, como la Ley de Ohm, la Ley de Faraday y la Ley de Lenz, y su aplicación en el diseño y análisis de motores y generadores eléctricos.

\subsection{Ley de Ohm}

La Ley de Ohm es esencial para el análisis de circuitos eléctricos. Establece que la diferencia de potencial\footnote{El cambio de energía potencial de una carga \(q\) desplazada de un punto \(A\) a un punto \(B\), dividido entre la carga.} (\(V\)) entre los extremos de un conductor es directamente proporcional a la corriente eléctrica (\(I\)) que circula por él, siendo la constante de proporcionalidad la resistencia (\(R\)) del conductor.

\begin{equation}
V = I \times R
\end{equation}

Donde:
\begin{itemize}
    \item \(V\): Voltaje aplicado [V]
    \item \(I\): Corriente eléctrica [A]
    \item \(R\): Resistencia [\Omega]
\end{itemize}

\ejemplo{Ejemplo 1}{Calcular la resistencia de un conductor por el que circula una corriente de 3 A cuando se aplica una tensión de 15 V.

Datos:
\begin{itemize}
    \item \(I = 3 \text{ A}\)
    \item \(V = 15 \text{ V}\)
\end{itemize}

Aplicamos la Ley de Ohm:
\begin{equation}
R = \frac{V}{I} = \frac{15 \text{ V}}{3 \text{ A}} = 5 \Omega
\end{equation}
}

\subsection{Ley de Faraday}

La Ley de Faraday de la inducción electromagnética establece que un cambio en el flujo magnético a través de un circuito induce una fuerza electromotriz\footnote{La fuerza electromotriz (fem) es el trabajo realizado por una fuente de energía por unidad de carga para mover cargas eléctricas a través de un circuito.} (\(\mathcal{E}\)) en el circuito.

\begin{equation}
\mathcal{E} = -\frac{d\Phi_B}{dt}
\end{equation}

Donde:
\begin{itemize}
    \item \(\mathcal{E}\): Fuerza electromotriz inducida [V]
    \item \(\Phi_B\): Flujo magnético [Wb]
\end{itemize}

El signo negativo indica que la fem inducida siempre se opone al cambio en el flujo magnético (Ley de Lenz).

\subsection{Ley de Lenz}

La Ley de Lenz complementa la Ley de Faraday y establece que la dirección de la corriente inducida en un circuito cerrado es tal que el campo magnético que crea se opone al cambio en el flujo magnético que lo produce. Esto se representa en la fórmula de la Ley de Faraday con el signo negativo.

\subsection{Aplicaciones de las Leyes Electromagnéticas en Máquinas Eléctricas}

Las máquinas eléctricas, como los motores y generadores, funcionan gracias a estas leyes electromagnéticas. Por ejemplo, en un generador, la rotación de una bobina en un campo magnético cambia el flujo magnético a través de la bobina, induciendo una fem y, por lo tanto, una corriente eléctrica.

\info{Fórmulas Clave}{
\begin{itemize}
    \item Ley de Ohm: \(V = I \times R\)
    \item Ley de Faraday: \(\mathcal{E} = -\frac{d\Phi_B}{dt}\)
\end{itemize}
}

\actividad{Responder el siguiente cuestionario}
\begin{enumerate}
    \item ¿Qué es la resistencia eléctrica y cómo se mide?
    \item ¿Cuál es la relación entre el flujo magnético y la fuerza electromotriz inducida según la Ley de Faraday?
    \item ¿Qué establece la Ley de Lenz respecto a la corriente inducida?
    \item ¿Cómo se aplica la Ley de Ohm en el análisis de circuitos eléctricos?
    \item ¿Qué unidades se utilizan para medir el flujo magnético y la fuerza electromotriz?
    \item ¿Cómo influye la velocidad de cambio del flujo magnético en la fuerza electromotriz inducida?
    \item ¿Qué papel juegan las leyes de Faraday y Lenz en el funcionamiento de un generador eléctrico?
    \item ¿Cómo se calcula la resistencia de un conductor si conocemos el voltaje aplicado y la corriente que circula?
    \item ¿Qué aplicaciones prácticas tienen las leyes electromagnéticas en la ingeniería eléctrica?
    \item ¿Cómo se relacionan las leyes de Faraday y Lenz en el contexto de un motor eléctrico?
\end{enumerate}

\actividad{Resolver los siguientes problemas}
\begin{enumerate}
    \item Calcular la corriente que circula por una resistencia de \(R = 10 \Omega\) cuando se aplica un voltaje de \(V = 20 \text{ V}\).
    \item Determinar la fuerza electromotriz inducida en una bobina cuando el flujo magnético cambia a razón de \(2 \text{ Wb/s}\).
    \item Calcular el voltaje necesario para que una corriente de \(5 \text{ A}\) circule por una resistencia de \(8 \Omega\).
    \item Hallar la resistencia de un conductor si una corriente de \(4 \text{ A}\) genera un voltaje de \(16 \text{ V}\).
    \item Determinar la corriente inducida en un circuito donde el flujo magnético disminuye a una tasa de \(0.1 \text{ Wb/s}\) y la resistencia es \(2 \Omega\).
    \item Calcular el cambio en el flujo magnético si se induce una fuerza electromotriz de \(3 \text{ V}\) en un tiempo de \(0.5 \text{ s}\).
    \item Hallar la resistencia si se sabe que una corriente de \(1 \text{ A}\) genera un voltaje de \(10 \text{ V}\).
    \item Determinar la corriente que circula por una resistencia de \(15 \Omega\) con un voltaje de \(30 \text{ V}\).
    \item Calcular la fuerza electromotriz inducida cuando el flujo magnético cambia en \(4 \text{ Wb}\) en \(2 \text{ s}\).
    \item Hallar la corriente si se conoce que una resistencia de \(5 \Omega\) genera un voltaje de \(25 \text{ V}\).
\end{enumerate}



\section{Ejemplos y Problemas Resueltos}

Para entender mejor las leyes electromagnéticas y su aplicación en máquinas eléctricas, a continuación se presentan ejemplos prácticos y problemas resueltos.

\subsection{Ejemplo de Ley de Faraday}

Supongamos que tenemos una bobina con 50 espiras en la que el flujo magnético cambia a razón de \(0.1 \, \text{Wb/s}\). Queremos calcular la fuerza electromotriz inducida en la bobina.

Datos:
\begin{itemize}
    \item \(N = 50\) (número de espiras)
    \item \(\frac{d\Phi_B}{dt} = 0.1 \, \text{Wb/s}\)
\end{itemize}

Aplicando la Ley de Faraday:
\begin{equation}
\mathcal{E} = -N \frac{d\Phi_B}{dt} = -50 \times 0.1 = -5 \, \text{V}
\end{equation}

Por lo tanto, la fuerza electromotriz inducida es \(-5 \, \text{V}\).

\section{Cuestionario}

\actividad{Responder el siguiente cuestionario}
\begin{enumerate}
    \item ¿Qué es la resistencia eléctrica y cómo se mide?
    \item ¿Cuál es la relación entre el flujo magnético y la fuerza electromotriz inducida según la Ley de Faraday?
    \item ¿Qué establece la Ley de Lenz respecto a la corriente inducida?
    \item ¿Cómo se aplica la Ley de Ohm en el análisis de circuitos eléctricos?
    \item ¿Qué unidades se utilizan para medir el flujo magnético y la fuerza electromotriz?
    \item ¿Cómo influye la velocidad de cambio del flujo magnético en la fuerza electromotriz inducida?
    \item ¿Qué papel juegan las leyes de Faraday y Lenz en el funcionamiento de un generador eléctrico?
    \item ¿Cómo se calcula la resistencia de un conductor si conocemos el voltaje aplicado y la corriente que circula?
    \item ¿Qué aplicaciones prácticas tienen las leyes electromagnéticas en la ingeniería eléctrica?
    \item ¿Cómo se relacionan las leyes de Faraday y Lenz en el contexto de un motor eléctrico?
\end{enumerate}

\section{Problemas}

\actividad{Resolver los siguientes problemas}
\begin{enumerate}
    \item Calcular la corriente que circula por una resistencia de \(R = 10 \Omega\) cuando se aplica un voltaje de \(V = 20 \text{ V}\).
    \item Determinar la fuerza electromotriz inducida en una bobina cuando el flujo magnético cambia a razón de \(2 \text{ Wb/s}\).
    \item Calcular el voltaje necesario para que una corriente de \(5 \text{ A}\) circule por una resistencia de \(8 \Omega\).
    \item Hallar la resistencia de un conductor si una corriente de \(4 \text{ A}\) genera un voltaje de \(16 \text{ V}\).
    \item Determinar la corriente inducida en un circuito donde el flujo magnético disminuye a una tasa de \(0.1 \text{ Wb/s}\) y la resistencia es \(2 \Omega\).
    \item Calcular el cambio en el flujo magnético si se induce una fuerza electromotriz de \(3 \text{ V}\) en un tiempo de \(0.5 \text{ s}\).
    \item Hallar la resistencia si se sabe que una corriente de \(1 \text{ A}\) genera un voltaje de \(10 \text{ V}\).
    \item Determinar la corriente que circula por una resistencia de \(15 \Omega\) con un voltaje de \(30 \text{ V}\).
    \item Calcular la fuerza electromotriz inducida cuando el flujo magnético cambia en \(4 \text{ Wb}\) en \(2 \text{ s}\).
    \item Hallar la corriente si se conoce que una resistencia de \(5 \Omega\) genera un voltaje de \(25 \text{ V}\).
\end{enumerate}

\end{document}

